\chapter{Introdução}
\label{chap:introduction}

A transformação digital tem impactado progressivamente os serviços
de saúde no Brasil, impulsionando o desenvolvimento de soluções
tecnológicas que visam otimizar processos e melhorar a experiência
de pacientes e profissionais de saúde. No entanto, o agendamento
de consultas e atendimentos permanece predominantemente dependente
de canais tradicionais como telefone e aplicativos de mensagens.

Segundo \citeonline{doctoralia2022}, 91\% dos centros médicos
brasileiros utilizam o telefone como principal canal de
agendamento, sendo que cerca de 70\% recebem mais de 20 ligações
diárias de pacientes, das quais uma em cada três não é atendida.
Esse cenário aponta para a demanda por soluções digitais que
modernizem e democratizem o acesso ao agendamento de serviços
de saúde no país.

\section{Problema}
\label{sec:problem}

A ausência de plataformas acessíveis e integradas de agendamento
de serviços de saúde impõe dificuldades tanto aos pacientes
quanto aos profissionais e clínicas. Do ponto de vista do
paciente, a dependência de canais síncronos limita o acesso
ao agendamento fora do horário comercial e aumenta o tempo
necessário para confirmar um atendimento. Para os profissionais
e clínicas, a ausência de um sistema centralizado dificulta a
visualização da agenda, aumenta a ocorrência de faltas não
comunicadas e sobrecarrega a equipe administrativa com tarefas
manuais e repetitivas.

Essa realidade não se restringe às clínicas médicas
convencionais. Profissionais de diversas áreas da saúde, como
psicólogos, psiquiatras, dentistas e nutricionistas, enfrentam
os mesmos desafios de gestão de agenda, uma vez que também
dependem de canais informais para organizar seus atendimentos.

Além disso, muitos pacientes não sabem qual especialidade
procurar diante de seus sintomas, o que resulta em consultas
com profissionais inadequados para o caso, gerando retrabalho
para o paciente e ocupação desnecessária da agenda dos
profissionais. Essa lacuna de orientação representa um problema
adicional que plataformas tradicionais de agendamento não
endereçam.

\section{Justificativa}
\label{sec:justification}

O desenvolvimento de uma plataforma de agendamento de serviços
de saúde representa uma oportunidade de aplicar conceitos
fundamentais da engenharia de software, tais como arquitetura
de sistemas, modelagem de dados, autenticação, infraestrutura
em nuvem e inteligência artificial, em um contexto relevante
tanto do ponto de vista técnico quanto social.

O mercado brasileiro de healthtechs tem crescido de forma
acelerada. Segundo \citeonline{abstartupsdeloitte2022}, o número
de startups de saúde no Brasil quase duplicou entre 2019 e 2021.
Esse crescimento se manteve nos anos seguintes: de acordo com
\citeonline{distrito2024}, o Brasil concentra 64,8\% das
healthtechs da América Latina, com crescimento de 37,6\% nos
investimentos em 2024. No segmento de agendamentos
especificamente, \citeonline{ligaventures2024} identificou
buscadores e agendamentos como uma das principais categorias
de healthtechs ativas no país, evidenciando a demanda por
soluções digitais nessa área.

Diferentemente de plataformas existentes que focam
predominantemente em agendamentos médicos, a Medicano propõe
uma abordagem abrangente, contemplando diferentes áreas da
saúde e incorporando um chatbot de triagem inteligente baseado
em Large Language Models (LLMs). Esse componente permite que
o paciente descreva seus sintomas em linguagem natural e seja
direcionado automaticamente para a especialidade mais adequada,
reduzindo o risco de consultas inadequadas e melhorando a
eficiência do processo de agendamento.

O presente trabalho se justifica, portanto, tanto pela
relevância social do problema abordado quanto pela oportunidade
de consolidar competências técnicas em engenharia de software
e inteligência artificial aplicadas a um domínio de crescente
importância no ecossistema de inovação brasileiro.

\section{Objetivos}
\label{sec:objectives}

\subsection{Objetivo Geral}
\label{subsec:general-objective}

Desenvolver uma plataforma web de agendamento de serviços de
saúde que conecte pacientes a clínicas e profissionais de
diferentes especialidades, incorporando um chatbot de triagem
inteligente baseado em Large Language Models para orientar
o paciente à especialidade mais adequada às suas necessidades.

\subsection{Objetivos Específicos}
\label{subsec:specific-objectives}

\begin{itemize}
  \item Implementar sistema de autenticação para
        perfis de paciente e profissional de saúde

  \item Implementar gerenciamento de ciclo de vida
        de tokens JWT com Redis, permitindo revogação
        imediata de sessões no logout        ← aqui

  \item Desenvolver chatbot de triagem inteligente
        utilizando Vercel AI SDK integrado a modelos
        de linguagem de grande escala para
        direcionamento por especialidade

  \item Desenvolver módulo de busca de profissionais
        e clínicas por especialidade e localização

  \item Construir sistema de gerenciamento de slots
        de disponibilidade para profissionais e
        clínicas

  \item Implementar fluxo completo de agendamento,
        confirmação e cancelamento de atendimentos

  \item Configurar proxy reverso com Nginx e
        certificado SSL para exposição segura da
        API em produção

  \item Implementar documentação automática da API
        utilizando o padrão OpenAPI com Swagger

  \item Realizar o deploy do frontend na Amazon S3
        com distribuição via CloudFront e do backend
        em instância Amazon EC2
\end{itemize}

\section{Organização do Trabalho}
\label{sec:organization}

O presente trabalho está organizado em oito capítulos. O
Capítulo~\ref{chap:theoretical-background} apresenta o
referencial teórico que embasou as decisões técnicas do projeto.
O Capítulo~\ref{chap:methodology} descreve a metodologia de
desenvolvimento adotada. O Capítulo~\ref{chap:specification}
detalha a especificação do sistema, incluindo requisitos,
casos de uso e modelagem de dados. O
Capítulo~\ref{chap:architecture} documenta a arquitetura
da solução. O Capítulo~\ref{chap:schedule} apresenta o
cronograma de desenvolvimento. O Capítulo~\ref{chap:risks}
analisa os riscos identificados. Por fim, o
Capítulo~\ref{chap:conclusion} apresenta as conclusões e
trabalhos futuros.